\chapter{Literature Review and Research Gap}
\label{ch:literature}

\section{Rotary Inverted Pendulum as a Control Benchmark}
Furuta, Yamakita, and Kobayashi established the rotary inverted pendulum as a
compact benchmark for underactuated control, including the combined problems
of swing-up and stabilization \cite{Furuta1991}. Subsequent work has used the
platform to study linear control, nonlinear control, intelligent control,
parameter identification, and real-time education. The review by Hamza
et al.\ summarizes this development and shows that the platform remains useful
because a simple mechanical arrangement exposes nonlinear coupling,
open-loop instability, actuator constraints, and implementation sensitivity
\cite{Hamza2019}.

The present work uses the same benchmark for a narrower purpose: local
upright stabilization and base-reference motion on a custom stepper-driven
rig. Swing-up is excluded so that the experimental comparison focuses on
upright controller structure and actuator realization.

\section{Dynamic Modelling and Parameter Evaluation}
Furuta pendulum models are commonly derived with Euler--Lagrange equations
using motor torque as the generalized input. Cazzolato and Prime show that
apparently small choices in coordinates, inertia definitions, and mass
distribution can change the resulting equations, and provide a careful
reference derivation for the coupled dynamics \cite{Cazzolato2011}. Their work
supports the modelling structure used in this thesis.

For a custom apparatus, geometric measurement alone does not guarantee an
accurate dynamic model. Garc\'ia-Alarc\'on et al.\ discuss parameter
identification for a Furuta pendulum and demonstrate how physical parameters
can be estimated from experimental data \cite{GarciaAlarcon2012}. The present
thesis instead uses measured geometry and mass properties to obtain a
design-oriented model. This is sufficient for controller construction, but the
resulting constants are treated as nominal model values rather than a
high-fidelity identified plant.

Most standard models retain torque or motor voltage as the plant input. In the
present apparatus, however, the low-level interface commands step rate. The
model is therefore reduced to an abstract commanded base acceleration for
controller design. This choice improves consistency between equations and
firmware, but does not imply that the physical stepper realizes the requested
acceleration.

\section{Linear Upright Stabilization}
Linearization about the upright equilibrium enables state-space control with
pole placement or linear-quadratic regulation. Such controllers are widely
used as baselines because they provide transparent gain interpretation and
are computationally inexpensive. Park, Kim, and Lee combine swing-up with LQR
stabilization, while Chawla and Singla report a robust LQR-based intelligent
controller on a real rotary pendulum \cite{Park2011,Chawla2018}. Standard
state-feedback and pole-placement principles are described by Ogata and by
{\AA}str{\"o}m and Murray \cite{ogata_modern_control,Astrom2008}.

Linear controllers are locally meaningful: their design model is accurate
only near upright and within the actuator range. Their practical value remains
high because they provide a stable reference against which nonlinear methods
can be assessed. In this thesis, the linear full-state controller is therefore
the baseline and is also used for the commanded-reference experiment.

\section{Sliding-Mode Control}
Sliding-mode control constructs a surface in the state space and uses a
reaching law to drive the system toward that surface. Classical treatments by
Utkin and by Slotine and Li establish the robustness and Lyapunov arguments
used in controller design \cite{utkin1992,slotine1991}. Young, Utkin, and
\"{O}zg\"{u}ner emphasize both the practical strengths of sliding modes and the
implementation problem of high-frequency switching or chattering
\cite{Young1999}. Boundary layers and saturation functions are common ways to
smooth the discontinuous sign term.

Several studies apply sliding-mode variants to rotary inverted pendulums.
Yi\u{g}it reports model-free sliding-mode stabilization on a real rotary
pendulum \cite{Yigit2017}. Nguyen et al.\ use a fuzzy super-twisting approach,
and Pham, Dao, and Nguyen present a hierarchical sliding-mode design for
swing-up and stabilization \cite{Nguyen2020,Pham2024}. Irfan et al.\ compare
advanced sliding-mode techniques in modelling and simulation
\cite{Irfan2018}. These studies motivate nonlinear control, but they also show
that performance depends strongly on the surface definition, reaching law,
state estimation, and actuator assumptions.

This thesis studies two practical variants. The first uses a pendulum-only
surface plus a separately gated base-centering term. The second places
pendulum and base-tracking states in one four-state surface. In both cases, the
ideal Lyapunov result is separated explicitly from the sampled, saturated,
stepper-driven implementation.

\section{Experimental and Educational Implementations}
Integrated pendulum platforms and educational rigs demonstrate the importance
of sensing, timing, and safe experiment workflows in addition to control-law
design. The STMicroelectronics educational platform documents a complete
rotary-pendulum implementation, and Ju, Lee, and Lee describe a recent
real-time educational platform for control and reinforcement-learning
experiments \cite{ST_RIP_2019,Ju2025}. These systems support repeatable
experiments through integrated actuation and instrumentation.

The custom platform in this thesis differs in two important respects. First,
both angles are measured directly, but the actuator remains an open-loop
STEP/DIR stepper axis rather than a torque-controlled servo. Second, the
analysis is built around raw session logs and a pinned manifest, allowing each
reported result to be traced back to a specific trial.

\section{Research Gap}
The literature provides mature Furuta dynamics, linear stabilization methods,
sliding-mode theory, and many specialized controller designs. The narrower
engineering gap addressed in this thesis is the documented implementation and
like-for-like qualitative evaluation of three upright controllers on the same
measured-axis, stepper-driven platform.

The research gap is not a claim that linear or sliding-mode control is new.
Rather, the contribution lies in connecting:
\begin{enumerate}
  \item a first-principles model and an acceleration-command representation;
  \item three controller structures using the same measured states and
  firmware interface;
  \item explicit actuator, sensor, timing, and safety constraints; and
  \item an auditable experimental analysis that distinguishes ideal-theory
  properties from observed hardware behavior.
\end{enumerate}

This framing leads directly to the platform and model developed in
Chapter~\ref{ch:platform_modelling}.
